We present \zeroshotmid\ \method\ in this section. To solve a task, \zeroshotmid\ \method\ employs an agent to create one set of instructions per dataset which allows an LLM to reason toward the final prediction. Intuitively, humans often rely on specific instructions to more effectively guide their thought process as they work towards a solution to a problem. For instance, to understand the sentiment in the movie reviews, instructions such as ``\textit{1. Understand the Dataset: ... Movie Reviews dataset ... 2. Analyze the Passage: Pay attention to ... the tone of the review ...}'' help humans decompose the problem into task-specific reasoning steps and solve each to deliver the final answer (Figure~\ref{fig:full-pipeline}). \zeroshotstart\ \method\ follows this intuition.
\begin{figure}
  \centering
  \includegraphics[width=\linewidth]{figures/pdfs/4-21-fig3-langchain-agent.drawio.png}
  \caption{An example of our agent producing task-specific instructions (highlighted) for a classification dataset IMDB. The agent only runs once to produce the instructions. Then, the instructions are used for all our models during reasoning. Agent instructions for all datasets are in Appendix~\ref{sec:all_examples}.
  }
  \label{fig:agent-details}
\end{figure}

\paragraph{Agent Instructions} Instead of handwriting task-specific instructions, we build an agent to automate the process. Our agent only needs to generate instructions once per task instead of running the agent on all dataset instances. The intuition is that an agent is able to synthesize high-quality instructions with access to a wide range of existing task knowledge on the web. We design our agent based on ReAct~\citep{yao2023react} using LangChain's zero-shot ReAct implementation~\citep{langchain}. This is motivated by the recent developments of language agents for task solving. Our agent highlights two features (Figure~\ref{fig:agent-details}): (\expandafter{\romannumeral1}) Instruction generation. Our agent follows ReAct which uses an LLM to propose a series of thoughts. The agent then receives observations and takes actions following the thoughts. Different from ReAct which aims to directly solve the task, our agent outputs step-by-step instructions on how to solve the task. The major advantage of this is that we only need to generate instructions once per dataset instead of running the agent on all dataset instances. We use GPT-4~\citep{openai2023gpt4} as the default agent. Once the action is \texttt{finish}, the corresponding output is our task-specific instructions. (\expandafter{\romannumeral2}) Action space.
We constrain our action space to contain two types of actions that support the instruction generation: 
(a) \texttt{ask\_about\_dataset}[string], which returns information from the top relevant web pages containing information about the dataset. To do this, we construct and utilize a question answering API as a tool. This API is meant to answer questions about the document by interfacing with Chroma~\citep{chroma}, a vector database storing the web pages, to provide an answer.
(b) \texttt{finish}[instructions], which finishes the instruction generation with the task-specific instructions. As shown in Figure~\ref{fig:agent-details}, to produce the instructions, our agent takes basic dataset information such as the name of the dataset (e.g., IMDB), a few input-only examples (examples without ground truth labels), and the set of output labels for the dataset (if applicable, else the type of dataset, e.g., generation) as its input. Using the task knowledge from the web, our agent forms observations (e.g., ``\textit{Observation 1: ... labeled as either positive or negative ...}'') and thoughts (e.g., ``\textit{Thought 2: ... creating instructions ...}'') which trigger the agent to perform actions, such as the \texttt{finish} action to output the task-specific instructions. For a full description of the \method\ pipeline for instruction generation, see Appendix~\ref{sec:implementation-details}.

\paragraph{Chain of Thought Reasoning} Chain of thought (CoT)~\citep{wei2023chainofthought, kojima2023large} prompts LLMs to break down the task into intermediate reasoning steps that lead to the final answer. Unlike zero-shot CoT which uses a fixed prompt ``Let's think step by step'', we prepend our task-specific agent instructions (created only once per dataset) to the input of whichever models we are using for reasoning. These same instructions can be used for various reasoning LLMs on all dataset instances. The instructions prompt the LLMs to optimize their reasoning processes for the task. The LLMs will then follow our task-specific instructions to decompose the task into a chain of more specific intermediate steps to solve the task. As shown in Figure~\ref{fig:full-pipeline}, the agent instructions ``\textit{... Pay attention to ... explicit or implicit expressions of sentiment towards the movie ...}'' are the key to producing the critical reasoning path ``\textit{... the movie is worth a view only for the performances of the three actors ...}'', which leads to the correct prediction where standard zero-shot and zero-shot CoT fail. We follow zero-shot CoT, which consists of a reasoning extraction prompt to produce the intermediate reasoning steps, and an answer extraction prompt to collect the answers. For simplicity of implementation, we replace zero-shot CoT's fixed prompt with zero-shot \method's task-specific instructions in the reasoning extraction prompt.

\zeroshotstart\ \method\ enjoys several unique properties: (\expandafter{\romannumeral1}) \zeroshotstart\ \method\ is a new way to improve the zero-shot reasoning of LLMs. \zeroshotstart\ \method\ decouples the language agent and reasoning process of LLMs, which helps \zeroshotmid\ \method\ generalize to more tasks. Importantly, our language agent only needs to run once to produce the instructions for each task. And, the agent instructions provide more task-specific controls to the reasoning paths of LLMs, which benefits human alignment and improves the safety of LLMs. (\expandafter{\romannumeral2}) Our agent instructions are customized for different tasks and verified by existing task knowledge. For each task, different LLMs use the same set of instructions. We find the instructions transfer well among these reasoning LLMs. This is important in practice as the agent LLMs are often more powerful and costly than the reasoning LLMs, so our approach is a cost-effective alternative to using agents directly. Following the teacher-student knowledge distillation setup, our work is focused on using a larger LLM-based agent to instruct the smaller reasoning LLM.
(\expandafter{\romannumeral3}) By providing task-specific instructions, chain of thought reasoning abilities are further generalized to more tasks beyond reasoning tasks. We show general language understanding tasks such as generation and classification also benefit from chain of thought reasoning. \method\ is zero-shot so no input-output examples are required to solve the task. To better utilize the emerging reasoning capabilities of LLMs, our task-specific instructions align the reasoning process with a particular task better than general or fixed instructions.